% ============================================================
% Chapter 6 – Discussion
% chapters/06-discussion.tex
% ============================================================

\chapter{Discussion}
\label{chap:discussion}

\todo[inline]{
\begin{itemize}
    \item - How it went: est ce que ca marche
    \item - Limitations. M. Debruyne disait "J'ai implementé X, mais Y nécessiterait plus temps à cause de telle contrainte technique"...
\end{itemize}
}

Critique quant au fait que, bien qu'il s'agisse maintenant d'un formulaire, donc une solution visuelle sans code, l'utilisateur est tout de même contraint à devoir connaitre le RML mapping pour pouvoir remplir ce formulaire RML. Et pour la génération d'autres formulaires, il doit pouvoir comprendre les shapes, ce qui n'est pas forcément trivial non plus. Donc, même si on a réussi à faire du no-code, il y a quand même une certaine expertise requise pour utiliser le formulaire de manière efficace.